\styledchapter{Cooperative Games}
\section{Introduction}

\begin{example} \label{4-29-2}
	Consider nearby cities A, B, and C, each of which needs to build a large public project, e.g. an airport or water treatment facility. They could cooperate or build the project individually.

	Define the costs for various cities building it individually or jointly as
	\begin{align*}
		c(A) &= 120, \quad c(B) = 140, \quad c(C) = 120 \\
		c(AB) &= 170, \quad c(AC) = 160, \quad c(BC) = 190 \\
		c(ABC) &= 255
	.\end{align*}

	Clearly, since $c(AB) < c(A) + c(B)$, so the planner would prefer cities A and B to work together than to work separately. How much should she charge each city to make up the total cost of $170$? A naive approach is to make each city pay $85$. Any payment rule in which both cities pay less than they would have individually is sensible. We might also think that city B should pay more because it has a higher individual cost --- we can take a savings-based approach, splitting the total savings of $260 - 170 = 90$ evenly, so A pays 75 and B pays 95.

	What if we consider a three-way cooperation? Then $c(ABC) = 255$ must satisfy four constraints:
	\begin{align*}
		c(ABC) &< c(A) + c(B) + c(C) \\
		c(ABC) &< c(A) + c(BC) \\
		c(ABC) &< c(AB) + c(C) \\
		c(ABC) &< c(AC) + c(B)
	,\end{align*}
	all of which indeed hold. So the cheapest way to implement the project is to fund it together.

	Now, how should all three cities split the cost? Doing so evenly leads to everyone paying $85$, but then A and C are collectively paying more than $c(AC)$. The savings-based approach leads to everyone getting savings of $\frac{125}{3}$ over their individual costs. In this case, A pays $78 \frac{1}{3}$, B pays $98 \frac{1}{3}$, and C pays $78 \frac{1}{3}$ --- now A and B are collectively paying more than $c(AB) = 170$!
\end{example}

We just had three agents, but the number of comparisons we must make is dizzying.
Let's develop a framework for characterizing solutions in which agents are happy at the three-way joint outcome.

Let $x_A, x_B, x_C$ be the amounts paid by A,B, and C respectively.
Then we require the \emph{no subsidy constraints}:\boxednote{We can define the no subsidy constraints in a general situation by replacing the numbers with the relevant costs $c(A), c(AB)$, and so on.}
\begin{align*}
	x_A &\le 120, \quad x_B \le 140, x_C \le 120, \\
	x_A + x_B &\le 170, \quad x_A + x_C \le 160, \quad x_B+ x_C \le 190, \\
	x_A+x_B + x_C &= 255
.\end{align*}

Define the savings \[
	y_A \coloneqq c(A) - x_A, \quad y_B \coloneqq c(B) - x_B, \quad y_C \coloneqq c(C) - x_C
.\] 

We can define constraints in terms of the savings as\boxednote{Now, the numbers are differences of costs: $90 = c(A) + c(B) - c(AB)$.}
\begin{align*}
	y_A &\ge 0, \quad y_B \ge 0, \quad y_C \ge 0,\\
	y_A + y_B &\ge 90, \quad y_B + y_C \ge 70, \quad y_A + y_C \ge 80, \\
	y_A + y_B + y_C &= 125
.\end{align*}
We can solve this system of linear solutions to get that one solution in the feasible set is \[
	\left( y_A,y_B,y_C \right) = \left( 55, 35, 35 \right)
.\] 

Is it always possible to find such a solution?
If we increase the three-way cost to $c(ABC)' = 265$, the new system has the same one-way and two-way savings restrictions, but a different three-way savings constraint:
\begin{align*}
	y_A &\ge 0, \quad y_B \ge 0, \quad y_C \ge 0,\\
	y_A + y_B &\ge 90, \quad y_B + y_C \ge 70, \quad y_A + y_C \ge 80, \\
	y_A + y_B + y_C &= 115
.\end{align*}
But now summing all three two-way constraints gives \[
	(y_A + y_B) + (y_B + y_C) + (y_A + y_C) \ge 90 + 70 + 80
,\] which yields \[
	2 y_A + 2y_B + 2y_C \ge 240
,\] or \[
	y_A + y_B + y_C \ge 120
,\] 
which is incompatible with the new constraint $y_A + y_B + y_C = 115$. So there is no three-way cooperative outcome that is stable, despite the three-way cooperative outcome being the best (cheapest).

\section{Cooperative Games with Transferable Utility}

\begin{definition}[Cooperative game with transferable utility]
	A cooperative game with transferable utility is a pair $(N,v)$, where $N$ is a set of agents and\boxednote{The game has \emph{transferable} utility because there is a single utility for the coalition, not a vector of individual utilities for each member of the coalition.}
	\[
		v: 2^N \setminus \O \to \R
	,\] where $v(S)$ is the surplus associated with every nonempty coalition $S$.
	\boxednote{In Example (\ref{4-29-2}), the value $v(S)$ was the total savings of coalition $S$.}
\end{definition}

\begin{definition}[Allocation]
	An allocation is a vector $x \in \R^n$ such that \[
		\sum\limits_{i \in N}^{} x_i = v(N)
	.\] 
	The allocation specifies the amount of surplus that each agents gets when the coalition includes all agents.
\end{definition}
\begin{definition}[Core]
	The core of $(N,v)$ is the set of all allocations such that the sum of the allocation components within the coalition is larger than the surplus generated from the coalition: \[
		\left\{\text{allocations } x \in \R^n : \text{for all nonempty } S \subsetneq N, \sum\limits_{i \in S}^{} x_i \ge v(S)\right\}
	.\] 
	Alternatively and more clearly, it is \[
		\left\{x \in \R^n: \text{for all nonempty } S \subsetneq N, \sum\limits_{i \in S}^{} x_i \ge v(S) \text{ and } \sum\limits_{i \in N}^{} x_i = v(N)\right\}
	.\] 
\end{definition}

Consider $N= \left\{1,2,3\right\}$. Then the core can be written as all $x \in \R^3$ such that \[
	\begin{pmatrix}
		1 & 0 & 0 \\
		0 & 1 & 0 \\
		0 & 0 & 1 \\
		1 & 1 & 0 \\
		1 & 0 & 1 \\
		0 & 1 & 1
	\end{pmatrix} \begin{pmatrix}
	x_1 \\
	x_2 \\
	x_3
	\end{pmatrix} 
	\ge \begin{pmatrix}
	v\left( \left\{1\right\} \right) \\
	v\left( \left\{2\right\} \right) \\ v\left( \left\{3\right\} \right) \\
	v\left( \left\{1,2\right\} \right) \\ v\left( \left\{1,3\right\} \right) \\ v\left( \left\{2,3\right\} \right)
	\end{pmatrix} , \quad \begin{pmatrix}
	1 &  1 & 1
	\end{pmatrix} \begin{pmatrix}
	x_1 \\
	x_2 \\ x_3
	\end{pmatrix} = \begin{pmatrix}
	v\left( \left\{1,2,3\right\} \right)
	\end{pmatrix} 
.\] 

Before moving to the general case, we need a new definition.
\begin{definition}[Balanced weights]
	A set of balanced weights is a function $\delta: 2^{N} \setminus \left\{\O, N\right\} \to [0,1]$ such that for every agent $i \in N$, \[
		\sum\limits_{S: i \in S}^{} \delta(S) = 1
	.\] 
\end{definition}
A set of balanced weights gives a weight to every nonsilly\footnote{Nonempty proper subset.} coalition such that when each agent sums the weights of the coalitions she is in, the sum is $1$.
While a usual coalition structure partitions agents so each agent is in exactly one coalition, a set of balanced weights allows each agent to be in multiple coalitions (excluding the empty and grand coalitions), so long as her coalition assignments have total weight one.
Furthermore, because $\delta(S)$ does not depend on $i$, given a coalition, everyone must have the same weight (spend the same time) in that coalition.

\begin{theorem}[Bonderava--Shapley]
	A cooperative game with transferable utility $(N,v)$ has a nonempty core if and only if for every set of balanced weights $\delta$, \[
		\sum\limits_{S \subsetneq N \mid S \ne \O}^{} \delta(S) v(S) \le v(N)
	.\] 
\end{theorem}

Bonderava--Shapley says that the core is nonempty iff for every artificial coalition arrangement, the total value generated by summing across all coalitions does not exceed the grand coalition's value.

\begin{proof}
	Fix an arbitrary $N$. 

	To write the constraints that define the core as a linear system, define $A \in \left\{0,1\right\}^{\left( 2^{\left| N \right|} - 2 \right) \times \left| N \right|}$ by \[
		A_{S,i} = \begin{cases}
			1, \quad & i \in S \\
			0, \quad &\text{otherwise}
		\end{cases}
	,\] 
	where $A$ has one row for each nonempty proper subset of $N$ (i.e. every possible nonsilly coalition) and one column for each agent. 
	Let $\vec{v} \in \R^{2^{\left| N \right|} - 2}$ have in its coordinate with index $S$ the surplus $v(S)$.

	Now we can characterize the core by\boxednote{Note that $\mathds{1}$ is a row vector. So later, $\mathds{1}'$ is a column vector.} \[
		\left\{x \in \R^n \mid Ax \ge \vec{v}, \mathds{1} \cdot x = v(N)\right\}
	.\] 

	Negating the inequality, Farkas's Lemma (\ref{4-1-2}) says that the core is nonempty iff there is no row matrix $p \ge 0$ and $q \in \R$ such that 
	\begin{align*}
		p\left( -A \right) + q \mathds{1}' &= 0 \\
		p(-\vec{v}) + qv(N) &< 0
	.\end{align*}
	If the system has a solution, then $q\ne 0$, as otherwise, we have $p=0$, which implies from the second equation that $0 (-\vec{v}) + 0 v(N) = 0 < 0$.
	Furthermore, $q \ge 0$ because each component of $p$ is nonnegative and each component of $A$ is either $0$ or $1$. Thus $q > 0$, and we can divide both equations by $q$ to get 
	\begin{align*}
		\overbrace{\left(\frac{1}{q}p\right)}^{\coloneqq \delta}A &= \mathds{1}' \\
		\left( \frac{1}{q}p \right)\vec{v} &> v(N) \\
		\left( \frac{1}{q}p \right) &\ge 0
	.\end{align*}
	From the first equation, we see that $\delta$ is a set of balanced weights, since for every agent $i \in N$ (corresponding to a column in $A$ and coordinate in the row vector $\mathds{1}'$), \[
		\sum\limits_{S: i \in S}^{} \delta(S) = 1
	.\] 
	The second equation translates to the inequality in the proposition statement, and the third requires becomes the definition restriction that weights are nonegative.

	We can think of $\delta$ as \emph{artificial coalitions} --- time allocations over actual coalitions.
\end{proof}

\section{Cooperative Games with Non-Transferable Utility}

Recall that in games with transferable utility, we had surplus within the coalition of $v(S)$ for every nonempty $S \subseteq N$.

In games with nontransferable utility, we assume that for every nonempty $S \subseteq N$, there is a set of actions $A_S$ that can be taken by the coalition $S$. Let $A \coloneqq \cup_S A_S$ be the set of all actions, and assume that each agent $i$ has a weak ranking $\succsim_i$ over her set of actions $\cup_{\left\{S : i \in S\right\}}A_S$.

\begin{definition}[Core]
	An action $a \in A_N$ is in the core if there is no coalition $S \subseteq N$ and $a' \in A_S$ such that \[
		a' \succsim_j a, \ \forall j \in S \text{ and } a' \succ_k a \text{ for some }k \in S
	.\] 
\end{definition}

Answers for when the core exists and on its properties are not very satisfying, so we will look at more specialized examples.
